\documentclass[9pt, mathserif]{beamer}

\usepackage[utf8]{luainputenc}

\usepackage[ngerman]{babel}
\setlength{\parindent}{0pt}
\usepackage[babel,german=quotes]{csquotes}
\usepackage{eurosym}

\usepackage{amsmath}
%\usepackage{mathtools}
\usepackage{amssymb}
\usepackage{xfrac}


\usepackage[locale=DE, separate-uncertainty=true, per-mode=fraction, alsoload=astro]{siunitx}

\usepackage{booktabs}
\usepackage{tabulary}
\usepackage{threeparttable}

\usepackage[]{graphicx}

\usepackage[]{hyperref}
\usepackage{braket}
\usepackage{xcolor}
\usepackage{enumerate}

\usepackage{listings}
\usepackage{lstautogobble}

\lstnewenvironment{lstverbatim}[1][]{
  \lstset
  {
    autogobble,
    language=TeX,                
    basicstyle=\ttfamily,
    commentstyle=\color{blue},
    columns=fullflexible,
    breaklines=true,
    keepspaces=true,
    aboveskip=-0.2em,
    belowskip=-0.2em
  }
}{}


\title{Verfassen wissenschaftlicher Texte mit \LaTeX}
\author{PeP et Al. toolbox}
\date{2014}
\institute[Pep et Al. e.V.]{\includegraphics[width=0.2\textwidth]{peplogo.pdf}}
\titlegraphic{./latex-vs-word.pdf}

\usetheme{TUDo}

\begin{document}

\begin{frame}
\setcounter{framenumber}{0}
    \titlepage
\end{frame}

\begin{frame}
    \frametitle{Übersicht}
    \tableofcontents
\end{frame}

\section{Zahlen und Einheiten}
\begin{frame}[fragile]
    \frametitle{Das SIunitx-Paket}
    Dieses Paket sollte immer und für jede Zahl mit oder ohne Einheit verwendet werden.

   \begin{block}{benötigte Pakete}
    \begin{lstverbatim}
        \usepackage[locale=DE, separate-uncertainty=true, per-mode=fraction]
                   {siunitx}
    \end{lstverbatim}
   \end{block}
    \begin{columns}[T]
        \begin{column}{0.6\textwidth}
            \begin{block}{\LaTeX-Code}
                \begin{lstverbatim}
                \num{1.23456} und \num{987654321}
                \num{1.2e2}
                \si{\newton} = \si{\kilo\gram\metre\per\second\squared}
                \SI{1.2}{\metre\per\second}
                \SI{4,3(12)}{\micro\second}
                \end{lstverbatim}
            \end{block}
        \end{column}
        \begin{column}{0.35\textwidth}
            \begin{block}{Ergebnis}
                \num{1.23456} und \num{987654321} \\
                \num{1.2e2} \\
                \si{\newton} = \si{\kilo\gram\metre\per\second\squared} \\
                \ \\
                \SI{1.2}{\metre\per\second} \\
                \SI{4,3(12)}{\micro\second}
            \end{block}
        \end{column}
    \end{columns}
\end{frame}

\section{Formelsatz}
\subsection{im Fließtext}
\begin{frame}[fragile]
   \frametitle{\$ \dots \$-Umgebung}
   Für mathematische Symbole, Variablen und kleine Formeln im Fließtext. 

   \begin{block}{benötigte Pakete}
    \begin{lstverbatim}
        \usepackage{amsmath}
        \usepackage{mathtools}
        \usepackage{amssymb}
    \end{lstverbatim}
   \end{block}
    \begin{columns}[T]
        \begin{column}{0.6\textwidth}
            \begin{block}{\LaTeX-Code}
                \begin{lstverbatim}
                $x$ 
                $x^i$
                $x^12$ bzw. $x^{12}$ % Vorsicht
                $x_\text{max}$
                $U(t) = U_0 \cdot \cos(\omega t)$
                \end{lstverbatim}
            \end{block}
        \end{column}
        \begin{column}{0.35\textwidth}
            \begin{block}{Ergebnis}
                $x$  \\
                $x^i$ \\
                $x^12$ bzw. $x^{12}$ \\ 
                $x_\text{max}$ \\
                $U(t) = U_0 \cdot \cos(\omega t)$
            \end{block}
        \end{column}
    \end{columns}
\end{frame}


\end{document}
